\section{Interaction-Presumption Strength and Changes in Causal Effects}
Research Question~\ref{rq:correlation_diff}, introduced in Section~\ref{sec:research_questions}, examines whether joint interventions have a greater effect when the oracle expresses a stronger presumption of interaction among the concepts within a question. Specifically, it asks whether the differences in question-averaged causal effects and question-level faithfulness between the joint and atomic formulations increase with interaction-presumption strength. As described in Section~\ref{sec:analysis_procedures}, this quantity is operationalized by $\Delta\kappa$, which is derived from the candidate interaction groups and their implied treatments relative to the number of concepts. It is a property of the oracle output rather than a direct measure of statistical association or response-level interaction.

Figures~\ref{fig:hyp1-bbq-scatterplots} and~\ref{fig:hyp1-medqa-scatterplots} show the relationship between interaction-presumption strength and the differences in causal effects and faithfulness. Each blue point represents one question, and the orange line is the least-squares linear fit. In the BBQ panels, the points exhibit no clear positive pattern; the estimated associations are weak. The MedQA panels suggest a possible positive association for $\Delta\mathrm{CE}$, particularly with \texttt{gpt-oss:120b}.

% RQ1 preset: fixed dimensions, x variable, solid fitted line, and test summary.
% #1 Pearson r, #2 confidence-interval lower bound, #3 plot-specific keys.
\newcommand{\HypOneScatter}[3]{%
  \ResultScatterPlot{x column=diff_correlation_strength,
    x label={$\Delta\kappa(\mathbf{x})$},
    line style=solid,
    minipage width=0.485\textwidth,
    plot width=\linewidth,
    plot height=5.45cm,
    mark size=1.65pt,
    below skip=-0.25em,
    text below={Pearson $r=#1$; one-sided 95\% CI $[#2,\,1.000]$},
    #3
  }%
}

\begin{figure}[p]
  \centering
  \HypOneScatter{-0.025}{-0.329}{title={GPT-OSS 120B --- $\Delta\mathrm{CE}$},
    data=data/rq1/bbq_gptoss.csv,
    y column=mean_absolute_ce_diff,
    y label={$\Delta\mathrm{CE}(\mathbf{x})$},
    line intercept=0.04781242,
    line slope=-0.005150228,
    xmax=2.3,
    ymax=0.26
  }
  \hfill
  \HypOneScatter{-0.102}{-0.396}{title={GPT-OSS 120B --- $\Delta\mathcal{F}$},
    data=data/rq1/bbq_gptoss.csv,
    y column=faithfulness_diff,
    y label={$\Delta\mathcal{F}(\mathbf{x})$},
    line intercept=0.108172,
    line slope=-0.04293996,
    xmax=2.3,
    ymax=0.45
  }

  \par\vspace{1.1em}
  \HypOneScatter{0.094}{-0.219}{title={Mistral Medium 3.5 --- $\Delta\mathrm{CE}$},
    data=data/rq1/bbq_mistral.csv,
    y column=mean_absolute_ce_diff,
    y label={$\Delta\mathrm{CE}(\mathbf{x})$},
    line intercept=0.03854958,
    line slope=0.01621849,
    xmax=2.3,
    ymax=0.26
  }
  \hfill
  \HypOneScatter{0.119}{-0.195}{title={Mistral Medium 3.5 --- $\Delta\mathcal{F}$},
    data=data/rq1/bbq_mistral.csv,
    y column=faithfulness_diff,
    y label={$\Delta\mathcal{F}(\mathbf{x})$},
    line intercept=0.1011828,
    line slope=0.03338781,
    xmax=2.3,
    ymax=0.36
  }

  \caption{Relationship between increased oracle-derived interaction-presumption strength and
  changes in causal effects or question-level faithfulness for BBQ.}
  \label{fig:hyp1-bbq-scatterplots}
\end{figure}

\begin{figure}[p]
  \centering
  \HypOneScatter{0.338}{0.030}{title={GPT-OSS 120B --- $\Delta\mathrm{CE}$},
    data=data/rq1/medqa_gptoss.csv,
    y column=mean_absolute_ce_diff,
    y label={$\Delta\mathrm{CE}(\mathbf{x})$},
    line intercept=-0.0007759744,
    line slope=0.001300082,
    xmax=8,
    ymax=0.025
  }
  \hfill
  \HypOneScatter{0.236}{-0.082}{title={GPT-OSS 120B --- $\Delta\mathcal{F}$},
    data=data/rq1/medqa_gptoss.csv,
    y column=faithfulness_diff,
    y label={$\Delta\mathcal{F}(\mathbf{x})$},
    line intercept=0.01957124,
    line slope=0.004748344,
    xmax=8,
    ymax=0.18
  }

  \par\vspace{1.1em}
  \HypOneScatter{0.097}{-0.256}{title={Mistral Medium 3.5 --- $\Delta\mathrm{CE}$},
    data=data/rq1/medqa_mistral.csv,
    y column=mean_absolute_ce_diff,
    y label={$\Delta\mathrm{CE}(\mathbf{x})$},
    line intercept=0.001771789,
    line slope=0.0001623979,
    xmax=8,
    ymax=0.015
  }
  \hfill
  \HypOneScatter{-0.075}{-0.409}{title={Mistral Medium 3.5 --- $\Delta\mathcal{F}$},
    data=data/rq1/medqa_mistral.csv,
    y column=faithfulness_diff,
    y label={$\Delta\mathcal{F}(\mathbf{x})$},
    line intercept=0.04718424,
    line slope=-0.001921023,
    xmax=8,
    ymax=0.18
  }

  \caption{Relationship between increased oracle-derived interaction-presumption strength and
  changes in causal effects or question-level faithfulness for MedQA.}
  \label{fig:hyp1-medqa-scatterplots}
\end{figure}

As expected from the graphs, the correlation tests in Box~\ref{hyp:hyp1-results} identify only one supported association: for MedQA evaluated with \texttt{gpt-oss:120b}, $\Delta\kappa$ is significantly positively correlated with the difference in question-averaged causal effects. None of the other dataset--model and outcome combinations meet the prespecified decision criterion. The results therefore do not provide consistent evidence that stronger oracle presumptions of interaction are associated with greater changes in causal effects or faithfulness when joint interventions replace atomic interventions.
\HypothesisTestResults[label=hyp:hyp1-results]
  {One-sided Pearson tests for \rqref{rq:correlation_diff}. A combination
  supports the hypothesis when $r\geq0.10$ and $p\leq0.10$.}
  {data/rq1/hypothesis_test_results.csv}
